% META: {"id": "TSPE-COMBI-EX-008", "chapitre": "TSPE-COMBINATOIRE", "type_objet": "exercice", "capacites": ["TSPE-COMBINATOIRE-C1"], "capacites_codes": ["C1"], "methodes": ["M1"], "parcours": 1, "competences": ["representer", "calculer"], "duree_min": 15, "mode_creation": "ex_nihilo", "sources_inspiration": [], "fichier_tex": "chapitres/TSPE-COMBINATOIRE/exercices/TSPE-COMBI-EX-008.tex", "corrige_tex": "chapitres/TSPE-COMBINATOIRE/corriges/TSPE-COMBI-CO-008.tex", "authoring": {"PEDAGOGICAL_GAP_ID": "GAP-889EB3571843", "CAPACITY_ID": "C1", "EXERCISE_ROLE": "DIRECT_APPLICATION", "DIFFICULTY": 1, "AUTHORING_REASON": "la capacite porte sur le CHOIX d'une representation parmi quatre, et le seul exercice rattache imposait l'arbre ; l'eleve n'avait donc jamais a construire un tableau ni un diagramme, ni a dire pourquoi l'un convient mieux que l'autre", "ORIGINALITY_PROVENANCE": "ex_nihilo"}, "status": "generated"}
% BEGIN-VERIFY
% from itertools import product
% from sympy import Rational
% lancers = list(product(range(1, 7), repeat=2))
% assert len(lancers) == 36
% somme_huit = [c for c in lancers if sum(c) == 8]
% assert len(somme_huit) == 5
% assert sorted(somme_huit) == [(2, 6), (3, 5), (4, 4), (5, 3), (6, 2)]
% assert Rational(len(somme_huit), len(lancers)) == Rational(5, 36)
% somme_sept = [c for c in lancers if sum(c) == 7]
% assert len(somme_sept) == 6
% effectif, anglais, espagnol, deux = 30, 18, 14, 7
% union = anglais + espagnol - deux
% assert union == 25
% assert effectif - union == 5
% assert union - deux == 18
% assert (anglais - deux) + (espagnol - deux) == 18
% assert (anglais - deux) + (espagnol - deux) + deux + (effectif - union) == effectif
% END-VERIFY
\begin{exercice}{TSPE-COMBI-EX-008}{1}{15}
Les deux situations suivantes sont indépendantes.

\medskip
\textbf{A.} On lance deux dés cubiques équilibrés, l'un rouge et l'un bleu, et on
s'intéresse à la somme des deux faces obtenues.
\begin{enumerate}
  \item Dresser le tableau à double entrée donnant la somme pour chacune des
    issues.
  \item Combien d'issues donnent une somme égale à $8$ ? Les énumérer.
  \item En déduire la probabilité d'obtenir une somme de $8$. Quelle somme est la
    plus probable, et pourquoi ?
\end{enumerate}

\medskip
\textbf{B.} Dans une classe de $30$ élèves, $18$ suivent l'option anglais
renforcé, $14$ l'option espagnol renforcé, et $7$ suivent les deux.
\begin{enumerate}
  \setcounter{enumi}{3}
  \item Représenter la situation par un diagramme et déterminer le nombre
    d'élèves qui suivent au moins une des deux options.
  \item En déduire le nombre d'élèves qui n'en suivent aucune, puis le nombre
    d'élèves qui en suivent exactement une.
  \item Pour chacune des deux situations, dire quelle représentation était
    adaptée, et expliquer pourquoi celle de l'autre situation aurait été
    maladroite ici.
\end{enumerate}
\end{exercice}
